\section{Degree-corrected SBM: the canonical spectral profile classifier}
\label{sec:dcsbm-proof}

This section proves the degree-corrected result in \cref{thm:dc}.
The reconstruction used here is the projection in the observed-degree
inner product.  It is an explicit decoder specification: the preceding
Euclidean reconstruction is not silently identified with it.
All likelihoods below are Bernoulli likelihoods.  The condition
\(np^2\le C_1\) controls their approximation at a bounded \emph{log-likelihood}
scale.

Write \(Z_{ia}=\ind{z_i=a}\), \(\Theta=\diag(\theta_1,\ldots,\theta_n)\),
\(B=pC\), and
\[
 \Omega=\Theta ZBZ^T\Theta,\qquad
 D_0=\diag(\Omega\one),\qquad
 D=\diag(D_1,\ldots,D_n),\quad D_i=\sum_jA_{ij}.
\]
The matrix \(\Omega\) includes its diagonal; the observed graph is hollow.
Thus \(\|\Omega-\E A\|_{\op}=O(p)\).  In particular, every raw noise
bound used below is unchanged up to its fixed constant by this completion.
Throughout this section a generic constant may depend on the fixed
parameter bounds in \cref{thm:dc}, but never on \(n\).  Constants depending
also on a fixed initialization-basin radius \(M\) are denoted \(C(M)\).

We use the common good event \(\mathcal G\) supplied by
\cref{lem:normalized-rows} and the probability calculation in
\cref{sec:probability}.  Its consequences needed here are
\begin{align}
 c d\le D_i\le C d,\qquad
 \|A-\Omega\|_{\op}\le C\sqrt d,\qquad
 \max_i\|U_{i\cdot}\|_2\le Cn^{-1/2},                    \label{dcsbm:good}\\
 \max_{a,b}\abs{[Z^T(A-\Omega)Z]_{ab}}\le Cn,           \label{dcsbm:block-event}
\end{align}
together with the corresponding normalized spectral and degree
consequences for the genuinely deleted graphs.  The retained eigenvectors
\(U\) are orthonormal, and their eigenvalues \(\mu_1,\ldots,\mu_K\)
are the \(K\) largest in absolute value.  The degree eigenvector is retained.
The common lemmas establish these consequences from the model and pay
their failure probabilities; they are not extra statistical assumptions.
We now prove all remaining deterministic and local probabilistic steps.

\subsection{The degree-metric reconstruction and its hollow correction}

Put \(S=D^{-1/2}AD^{-1/2}\), \(\Pi=UU^T\), and
\begin{equation}
 H=D^{-1/2}\Pi D^{1/2},\qquad
 R=AH=D^{1/2}U\diag(\mu_1,\ldots,\mu_K)U^TD^{1/2}.
 \label{dcsbm:canonical}
\end{equation}
The second formula evaluates \(R\) from retained spectral information.
Although \(H\) need not be symmetric, \(H^2=H\) and \(DH=H^TD\):
it is an orthogonal projection for
\(\langle x,y\rangle_D=x^TDy\).  Since the degree eigenvector
is proportional to \(D^{1/2}\one\), \(H\one=\one\) and
\(R\one=D\one\).

For \(h_i=R_{ii}\), define
\begin{equation}
 s_i=\frac{h_i}{1-\Pi_{ii}},\qquad
 R^\circ=R+\diag(s_1,\ldots,s_n)(H-I).
 \label{dcsbm:hollow}
\end{equation}
On \(\mathcal G\), the denominators are positive for all sufficiently
large \(n\).  Direct substitution gives
\begin{equation}
 R^\circ_{ii}=0,\qquad R^\circ\one=D\one,\qquad
 R^\circ Q=AQ\quad
 \text{if }\col(Q)\subseteq\col(D^{-1/2}U).
 \label{dcsbm:exact-identities}
\end{equation}
The last identity follows because \(HQ=Q\).  The hollow correction need
not be symmetric; a symmetric block parameter is fitted by the ordered
Bernoulli criterion, equivalently by symmetrizing its pair contributions.
The row profile scores always use the original rows of \(R^\circ\).

\begin{lemma}[Entry and absolute-row bounds]
\label{dcsbm:entry-lemma}
On \(\mathcal G\),
\begin{equation}
 \max_{i,j}|R^\circ_{ij}|\le Cp,\qquad
 \max_i\sum_j|R^\circ_{ij}|\le Cd,\qquad
 \max_i|s_i|\le Cp .
 \label{dcsbm:entry-bounds}
\end{equation}
Every action in \cref{dcsbm:canonical,dcsbm:hollow} is computable from
\((D,U,\mu)\), with no access to individual edges after encoding.
\end{lemma}

\begin{proof}
The normalized adjacency has operator norm at most one.  Consequently
\[
 |R_{ij}|
 \le \sqrt{D_iD_j}\,\|U_{i\cdot}\|_2\|U_{j\cdot}\|_2
 \le Cd/n=Cp.
\]
Moreover
\(
 |H_{ij}|
 \le\sqrt{D_j/D_i}\,\|U_{i\cdot}\|_2\|U_{j\cdot}\|_2
 \le C/n
\).
These estimates give \(|s_i|\le Cp\) and the entry bound for \(R^\circ\);
summing its entries in absolute value gives the row bound.
The displayed spectral formulas also prove the computational statement.
\end{proof}

\subsection{A finite deterministic spectral initializer}

The row-normalized population centers are particularly simple even when
the \(\theta_i\)'s vary within blocks.  To see this without replacing the
observed degrees by their expectations, put
\[
 M_D=D^{-1/2}\Omega D^{-1/2},\qquad
 G_D=Z^T\Theta D^{-1}\Theta Z.
\]
The diagonal entries of \(G_D\) are between constant multiples of \(n/d\).
The \(K\) nonzero eigenvalues of \(M_D\) have absolute values bounded
away from zero, by the lower singular-value bound on \(C\).
An orthonormal basis of its range has the form
\begin{equation}
 V=D^{-1/2}\Theta ZG_D^{-1/2}O
 \label{dcsbm:population-basis}
\end{equation}
for an orthogonal \(K\)-by-\(K\) matrix \(O\).  Each row has norm
comparable to \(n^{-1/2}\), and its normalization is exactly the
\(z_i\)-th row of \(O\).  These \(K\) unit centers have pairwise distance
\(\sqrt2\).

By \cref{dcsbm:good},
\(\|S-M_D\|_{\op}\le C/\sqrt d\).  The signal/bulk gap and
Davis--Kahan therefore give an aligned basis with
\(\|U-V\|_F^2\le C/d\).  The elementary normalization inequality
\[
 \left\|\frac{u}{\|u\|_2}-\frac{v}{\|v\|_2}\right\|_2
 \le \frac{2\|u-v\|_2}{\|v\|_2}
\]
holds when \(u\ne0\), and holds with a harmless change of constant
under any fixed unit-vector convention for \(u=0\).  For the normalized
rows \(x_i\) and true unit centers \(c_a\), it follows that
\begin{equation}
 \sum_{i=1}^n\|x_i-c_{z_i}\|_2^2\le Cn/d.
 \label{dcsbm:embedding-sse}
\end{equation}

The following initializer requires no numerical \(K\)-means certificate.
Starting with all rows remaining, repeat \(K\) times: choose a remaining
row whose closed ball of radius \(r=1/8\) contains the largest number
of remaining rows, break ties in a fixed order, record that row, and
remove all remaining rows at distance at most \(2r=1/4\) from it.
Finally assign every original row to its nearest recorded representative.
If a step is undefined, use the algorithm's fixed fallback partition.
The procedure uses \(O(Kn^2)\) distance comparisons in fixed dimension.

\begin{lemma}[Ball-count initialization]
\label{dcsbm:initializer}
Suppose every true block contains at least \(\gamma n\) rows and
the unit centers are separated by \(\sqrt2\).  Put
\(\epsilon=1/32\).  If
\(\sum_i\|x_i-c_{z_i}\|_2^2\le E\) and
\(m=E/\epsilon^2<\gamma n/3\), the procedure above selects exactly
one representative for each true center and, after one common
permutation of labels, makes at most \(m\) errors.  In particular,
\cref{dcsbm:embedding-sse} implies an initial error fraction at most
\(C/d\) on \(\mathcal G\).
\end{lemma}

\begin{proof}
Call a row good if it is within \(\epsilon\) of its true center.
There are at most \(m\) bad rows.  Inductively, all good rows of each
previously selected block have been removed and all good rows of each
unselected block remain.  A radius-\(r\) ball centered at a good row
of any unselected block contains every good row of that block, since
\(2\epsilon<r\).  Its count is at least \(\gamma n-m>m\).
Hence the maximizing ball contains a good row from an unselected block.
Its selected representative is within \(r+\epsilon=5/32\) of that
block's center.  Every good row of this block is within
\(r+2\epsilon=3/16<2r\) of the representative, while every good row
of another block is at least
\(\sqrt2-r-2\epsilon>2r\) away.  This proves the induction.
After \(K\) selections, a good row is at most \(3/16\) from its own
representative and at least \(\sqrt2-3/16\) from every other one.
Thus final nearest-representative assignment classifies all good
rows correctly.
\end{proof}

\subsection{Block-total accuracy and estimation of the nuisance parameters}

We first record the residual estimate behind both the block fit and the
tilted cancellation.

\begin{lemma}[Block residual in the degree metric]
\label{dcsbm:block-residual}
On \(\mathcal G\), with \(F=(I-\Pi)D^{1/2}Z\),
\begin{equation}
 \|F\|_{\op}\le\|F\|_F\le C\sqrt n,\qquad
 \sum_i\frac{(D_i-D_{0,i})^2}{D_i}\le Cn.
 \label{dcsbm:residual-estimates}
\end{equation}
The same conclusions hold for any genuinely deleted graph under its
corresponding common good event.
\end{lemma}

\begin{proof}
The degree-vector bound follows from
\[
 \|(D-D_0)\one\|_2
 =\|(A-\Omega)\one\|_2\le C\sqrt{nd}
\]
and \(\min_iD_i\ge cd\).
Write
\[
 D^{1/2}Z
 =D^{-1/2}D_0Z+D^{-1/2}(D-D_0)Z.
\]
Since \(D_{0,i}=p\theta_i\sum_c n_cC_{z_i c}\), the first term
belongs to \(\col(D^{-1/2}\Theta Z)=\col(M_D)\).
Its Frobenius norm is \(O(\sqrt{nd})\), while the angle between this
space and \(\col(U)\) is \(O(d^{-1/2})\).  Its residual is therefore
\(O(\sqrt n)\).  The squared Frobenius norm of the second term is
exactly the degree sum in \cref{dcsbm:residual-estimates}, because
every row of \(Z\) contains one unit entry.  This proves the assertion.
Deleting one vertex preserves all fixed group-mass bounds and the same
calculation applies to the minor.
\end{proof}

\begin{lemma}[Canonical reconstructed block totals]
\label{dcsbm:block-totals}
Fix \(M<\infty\).  Uniformly over every possibly data-dependent
partition \(\widehat Z\) with at most \(Mn/d\) mistakes after a common
permutation,
\begin{equation}
 [\widehat Z^TR^\circ\widehat Z]_{ab}
 =n_an_bpC_{ab}+O_M(n).
 \label{dcsbm:totals-accuracy}
\end{equation}
The usual diagonal exclusion is included in the \(O_M(n)\) term.
\end{lemma}

\begin{proof}
Let \(S_{\rm tail}=(I-\Pi)S(I-\Pi)\).  Spectral invariance gives the
exact quadratic identity
\begin{equation}
 Z^T(R-A)Z=-F^TS_{\rm tail}F .
 \label{dcsbm:quadratic-identity}
\end{equation}
Because \(M_D\) has rank \(K\), Weyl's inequality gives
\(\|S_{\rm tail}\|_{\op}\le C/\sqrt d\).
By \cref{dcsbm:block-residual}, the right side of
\cref{dcsbm:quadratic-identity} has norm at most \(Cn/\sqrt d\).
The hollow correction changes a true block total by at most \(Cd\):
indeed its total absolute contribution is at most
\[
 \max_i|s_i|\sum_i\sum_j|(H-I)_{ij}|
 \le Cp\,Cn=Cd.
\]
Together with \cref{dcsbm:block-event} and
\(Z^T\Omega Z=\diag(n_1,\ldots,n_K)(pC)\diag(n_1,\ldots,n_K)\),
this proves the claim for true labels.

Changing at most \(Mn/d\) labels changes at most \(C Mn^2/d\)
ordered pair memberships.  Each reconstructed entry has size at
most \(Cp\), so the additional error is at most \(C Mn\).
This deterministic argument is uniform over the entire basin and
requires no independence of \(\widehat Z\).
\end{proof}

For fitted labels \(\widehat z\), let \(\widehat n_a\) be their sizes,
\(\overline D_a^{\,\widehat z}\) their mean degrees, and set
\begin{equation}
 \widehat\theta_j=\frac{D_j}{\overline D_{\widehat z_j}^{\,\widehat z}}.
 \label{dcsbm:theta-fit}
\end{equation}
Thus \(\sum_{\widehat z_j=a}\widehat\theta_j=\widehat n_a\)
exactly.  Fit the symmetric matrix \(\widehat B\) by maximizing
\begin{equation}
 \frac12\sum_{i\ne j}\left\{
 R^\circ_{ij}\log(\widehat\theta_i\widehat\theta_j B_{\widehat z_i\widehat z_j})
 +(1-R^\circ_{ij})
       \log(1-\widehat\theta_i\widehat\theta_j B_{\widehat z_i\widehat z_j})
 \right\}
 \label{dcsbm:block-criterion}
\end{equation}
over positive symmetric \(B\), under a fixed probability cap below one.
This is a precisely defined signed Bernoulli working criterion.
The next lemma proves that its relevant roots exist and lie in the
interior on \(\mathcal G\).

\begin{lemma}[Uniform fitted-parameter bounds]
\label{dcsbm:nuisance}
For every fixed \(M\), uniformly over the basin in
\cref{dcsbm:block-totals}, the degree estimates are bounded above
and below by positive fixed constants and
\begin{align}
 \frac{\|\widehat\theta-\theta\|_2}{\sqrt n}
 &\le C(M)d^{-1/2},                                      \label{dcsbm:theta-rate}\\
 \max_{a,b}\left|\frac{\widehat B_{ab}}{pC_{ab}}-1\right|
 &\le C(M)(d^{-1}+p)\le C(M)d^{-1}.                       \label{dcsbm:B-rate}
\end{align}
The last inequality uses \(dp=np^2\le C_1\).
\end{lemma}

\begin{proof}
Every fitted group has size comparable to \(n\), and
\cref{dcsbm:good} bounds its mean degree above and below by
constant multiples of \(d\).  This proves boundedness of
\(\widehat\theta\).
With true labels, write \(D_i=dS_{z_i}\theta_i+\varepsilon_i\),
where \(S_a=\sum_c(n_c/n)C_{ac}\) and
\(\|\varepsilon\|_2\le C\sqrt{nd}\).
Division by the true-group mean degree, followed by subtraction
of \(\theta_i\), gives normalized squared error \(O(d^{-1})\).
Changing an \(\eta\)-fraction of labels changes each fitted-group
mean degree by at most \(Cd\eta\).  For correctly labeled vertices
this changes the estimated theta by \(O(\eta)\); incorrectly labeled
vertices are an \(\eta\)-fraction and have bounded error.  Hence
\[
 \frac{\|\widehat\theta-\theta\|_2}{\sqrt n}
 \le C(d^{-1/2}+\eta+\sqrt\eta),
\]
which yields \cref{dcsbm:theta-rate} for \(\eta\le M/d\).

For the block fit, group the ordered pairs in
\cref{dcsbm:block-criterion} belonging to the unordered block pair
\(\{a,b\}\).  Denote this ordered set by \(\mathcal J_{ab}\), put
\(\nu_{ab}=2-\ind{a=b}\), \(t_{ij}=\widehat\theta_i\widehat\theta_j\),
and \(E_{ab}^{\,\circ}=\sum_{\mathcal J_{ab}}R^\circ_{ij}\).
Its score equation is
\begin{equation}
 E_{ab}^{\,\circ}
 = b\sum_{(i,j)\in\mathcal J_{ab}}
      \frac{(1-R^\circ_{ij})t_{ij}}{1-bt_{ij}} .
 \label{dcsbm:block-root}
\end{equation}
All \(1-R^\circ_{ij}\) are positive for large \(n\).
The right side is strictly increasing in \(b\) below the cap.
By \cref{dcsbm:block-totals},
\[
 E_{ab}^{\,\circ}
 =\nu_{ab}n_an_bpC_{ab}+O_M(n)\asymp n^2p>0.
\]
For \(b\asymp p\), bounded theta and \(|R^\circ_{ij}|\le Cp\)
give
\[
 \sum_{\mathcal J_{ab}}
   \frac{(1-R^\circ_{ij})t_{ij}}{1-bt_{ij}}
 =\sum_{\mathcal J_{ab}}t_{ij}\{1+O(p)\}
 =\nu_{ab}\widehat n_a\widehat n_b
       \{1+O(p)+O(n^{-1})\}.
\]
The last equality uses exact fitted-block mean normalization;
for \(a=b\), removing the diagonal subtracts only
\(\sum_{\widehat z_i=a}\widehat\theta_i^2=O(n)\).
Evaluating \cref{dcsbm:block-root} at sufficiently small and large
fixed multiples of \(p\) brackets its unique root, so
\[
 \widehat B_{ab}
 =\frac{E_{ab}^{\,\circ}}{\nu_{ab}\widehat n_a\widehat n_b}
       \{1+O(p)+O(n^{-1})\}.
\]
Since \(|\widehat n_a-n_a|\le Mn/d\), this proves
\cref{dcsbm:B-rate}.  The root is \(O(p)\), hence satisfies the
fixed cap strictly for large \(n\).
\end{proof}

\subsection{Exact Bernoulli profiling at bounded score error}

Fix a tested vertex \(i\) with true block \(a\).  For candidate
block \(b\), define true neighbor masses and normalized profiles by
\begin{equation}
 M_c^{(i)}=\sum_{\substack{j\ne i\\z_j=c}}\theta_j,\qquad
 T_b^{(i)}=\sum_c M_c^{(i)}C_{bc},\qquad
 q_{bc}^{(i)}=\frac{M_c^{(i)}C_{bc}}{T_b^{(i)}} .
 \label{dcsbm:q-true}
\end{equation}
The fitted analogues are
\begin{equation}
 \widehat M_c^{(i)}
 =\sum_{\substack{j\ne i\\\widehat z_j=c}}\widehat\theta_j
 =\widehat n_c-\widehat\theta_i\ind{\widehat z_i=c},
 \qquad
 \widehat q_{bc}^{(i)}
 =\frac{\widehat M_c^{(i)}\widehat B_{bc}}
        {\sum_e\widehat M_e^{(i)}\widehat B_{be}}.
 \label{dcsbm:q-fitted}
\end{equation}
The factor \(p\) is not needed to evaluate the fitted quantity.
All these probabilities are bounded away from zero.
By \cref{dcsbm:nuisance} and fitted group-size accuracy,
\begin{equation}
 \max_{i,a,b,c}
 \left|\log\frac{\widehat q_{ac}^{(i)}}{\widehat q_{bc}^{(i)}}
          -\log\frac{q_{ac}^{(i)}}{q_{bc}^{(i)}}\right|
 \le C(M)d^{-1}.
 \label{dcsbm:q-rate}
\end{equation}
Indeed normalization and logarithms are smooth on the fixed compact
positive parameter set, and the relative mass error is
\(O_M(d^{-1}+n^{-1})\).

\begin{lemma}[Bernoulli profile reduction]
\label{dcsbm:profile-reduction}
Let a hollow row \(r=(r_j)_{j\ne i}\) have
\(\sum_jr_j=D_i\ge0\), \(\sum_j|r_j|\le Ld\), and \(D_i\le Ld\).
Suppose either \(r_j\in\{0,1\}\), or
\(\max_j|r_j|\le Lp=o(1)\).
For bounded positive neighbor parameters and block probabilities
of order \(p\), maximize the exact Bernoulli row criterion over
the target parameter \(t\ge0\), with a fixed probability cap below
one.  After dropping terms independent of the candidate label,
the profile contrast between \(a\) and \(b\) is
\begin{equation}
 L^{ab}_{r,\rm prof}
 =\sum_c N_c^r\log\frac{q_{ac}^{(i)}}{q_{bc}^{(i)}}
       +O(np^2),\qquad
 N_c^r=\sum_{\substack{j\ne i\\z_j=c}}r_j .
 \label{dcsbm:profile-linear}
\end{equation}
The same assertion holds with fitted labels, thetas, block probabilities,
and profiles.  Fixed positive candidate prior probabilities add their
bounded log-ratio to both sides.
\end{lemma}

\begin{proof}
For true parameters, omit the candidate-independent terms
\(\sum_jr_j\log\theta_j+D_i\log p\).  The remaining score is
\[
 \ell_r(b,t)
 =D_i\log t+\sum_cN_c^r\log C_{bc}
   +\sum_{j\ne i}(1-r_j)\log(1-tp\theta_jC_{bz_j}).
\]
For bounded \(t\), expansion of the logarithm gives uniformly
\begin{equation}
 \ell_r(b,t)
 =D_i\log t+\sum_cN_c^r\log C_{bc}-tpT_b^{(i)}
     +O(np^2).
 \label{dcsbm:profile-expansion}
\end{equation}
The restored edge part of the linear exposure is bounded by
\(Cp\sum_j|r_j|\le Cdp=Cnp^2\).
The quadratic remainder is bounded by
\(Cp^2(n+\sum_j|r_j|)\le Cnp^2\).

This expansion remains valid after maximizing.  In both allowed
row types \(1-r_j\ge0\), so the exact score is concave in \(t\).
For a raw row, at least \(n-O(d)\) nonedge terms remain; for a
reconstructed row, every nonedge coefficient is at least \(1/2\)
for large \(n\).  Thus the nonedge part of its derivative has
magnitude at least \(cd\), whereas the positive term is \(D_i/t\).
The derivative is negative above a sufficiently large fixed \(t_0\),
uniformly in the candidate.  The cap allows this bounded interval
for large \(n\), and the leading optimizer
\(t_b^0=D_i/(pT_b^{(i)})\) is also bounded.
Taking suprema in \cref{dcsbm:profile-expansion} is therefore valid.
The candidate-dependent leading maximum is
\(\sum_cN_c^r\log C_{bc}-D_i\log T_b^{(i)}\).
Since \(\sum_cN_c^r=D_i\), its difference is exactly the sum
in \cref{dcsbm:profile-linear}.

With fitted parameters replace \(pC\) by \(\widehat B\) and
the true masses by their fitted analogues.  Boundedness and
\(\widehat B\asymp p\) give the identical remainder calculation.
When \(D_i=0\), use the continuous extension after combining the
terms into \(D_i\log t\), with \(0\log0=0\).  Nonnegative nonedge
coefficients make \(t=0\) optimal, and the same contrast identity
holds.  In particular a raw zero row gives zero contrast.
\end{proof}

Exact preservation of \(D_i\) is essential in this argument:
it fixes the coefficient of \(\log t\) and the profiled total-degree
term.  The proof needs no \(O(d^{-1})\) estimate of an individual
target theta.

\begin{lemma}[Uniform replacement of all fitted stages]
\label{dcsbm:uniform-score}
Fix \(M\).  For every partition in the \(M/d\) basin, fitted by
\cref{dcsbm:theta-fit,dcsbm:block-criterion}, its actual row-profile
contrast satisfies, simultaneously for all \(i,a,b\),
\begin{equation}
 \widehat L_i^{ab}
 =[R^\circ w_{i,ab}]_i+O_M(1),\qquad
 (w_{i,ab})_j=\log\frac{q_{a,z_j}^{(i)}}{q_{b,z_j}^{(i)}} .
 \label{dcsbm:uniform-contrast}
\end{equation}
The formula is uniform over data-dependent labels and parameters
obtained by these fits, under the same common label permutation.
\end{lemma}

\begin{proof}
The profile remainder is \(O(np^2)=O(1)\).
Replacing \(\widehat Z\) by \(Z\) in a bounded block coefficient
vector changes row \(i\)'s score by at most
\[
 \max_j|R^\circ_{ij}|\,
     C\#\{j:\widehat z_j\ne z_j\}
 \le Cp\,CMn/d=C(M).
\]
Replacing the fitted coefficients by their true counterparts
costs at most
\(
 \sum_j|R^\circ_{ij}|\,
 C(M)d^{-1}\le C(M)
\)
by \cref{dcsbm:q-rate}.  These are deterministic bounds holding
for every partition in the basin, which proves the claimed uniformity.
\end{proof}

\subsection{The profile reference and its sharp directional tails}

For the fixed target \(i\), true block \(a\), and competitor \(b\),
write \(w_j=(w_{i,ab})_j\), \(v=\theta_i\), and
\[
 p_j=vp\theta_jC_{az_j},\qquad
 T_{i,ab}=\sum_{j\ne i}w_jA_{ij},\qquad
 K_{i,ab}(s)=\sum_{j\ne i}\log\{1+p_j(e^{-sw_j}-1)\}.
\]
Set
\begin{equation}
 J_{i,a\to b}
 =-\min_{s>0}K_{i,ab}(s).
 \label{dcsbm:profile-exponent}
\end{equation}
This is the exact Bernoulli transform of the leading profiled
contrast.  It is not asserted to be the simple-hypothesis
Chernoff information for the known-target-theta oracle.

\begin{lemma}[Uniform profile saddle and reference scale]
\label{dcsbm:profile-tail}
The minimizing \(s=s_{i,ab}\) stays in a fixed compact subinterval
of \((0,1)\).  Under the exact tilted product law,
\begin{equation}
 r_j=\frac{p_je^{-sw_j}}{1-p_j+p_je^{-sw_j}},\qquad
 \sum_j r_jw_j=0,\qquad
 \sum_j r_j(1-r_j)w_j^2\asymp d .
 \label{dcsbm:tilt}
\end{equation}
Uniformly for bounded real thresholds \(x\),
\begin{equation}
 \PP_a(T_{i,ab}<x)\asymp
 d^{-1/2}\exp\{-J_{i,a\to b}\}.
 \label{dcsbm:directional-tail}
\end{equation}
The exact full-edge Bernoulli profile-reference error at vertex \(i\)
is comparable, by fixed constants, to
\begin{equation}
 \sum_{b\ne a}d^{-1/2}\exp\{-J_{i,a\to b}\}.
 \label{dcsbm:reference-sum}
\end{equation}
\end{lemma}

\begin{proof}
Uniformly on bounded \(s\) intervals, the transform and its first
two derivatives satisfy
\[
 d^{-1}K_{i,ab}(s)
 =v\,\frac{T_a^{(i)}}n
    \left\{\sum_c(q_{ac}^{(i)})^{1-s}(q_{bc}^{(i)})^s-1\right\}
       +O(p).
\]
The derivative of the leading expression is strictly negative
at zero and strictly positive at one.  Positivity, the lower
singular-value bound on \(C\), and the group-mass bounds make
these signs and the positive second derivative uniform.
Indeed two normalized rows could coincide only if two rows of
\(C\) were proportional; compactness of the stated parameter
class excludes even an arbitrarily small normalized separation.
Strict convexity now gives the interior saddle and tilted
variance of order \(d\).

The bounded-weight, variance, and third-moment hypotheses of
\cref{lem:sharp} hold, so its bounded-shift tail conclusion proves
\cref{dcsbm:directional-tail}.  In particular, no nonlattice
assumption is necessary.  Alternatively, under the displayed
tilt, Berry--Esseen bounds every unit-interval mass by \(C/\sqrt d\);
summing against \(e^{sT}\) gives the upper bound.  A sufficiently
long fixed interval immediately below \(x\) has Gaussian mass
larger than twice the Berry--Esseen error, giving the lower bound.

The probability of \(D_i>Ld\) can be made
\(\exp(-Ld\,c_L)\) with arbitrarily large fixed exponential rate
by increasing the fixed upper cutoff.  On its complement,
\cref{dcsbm:profile-reduction} compares the exact reference
contrast to \(T_{i,ab}\) up to a fixed additive constant.
The upper-degree exception is negligible relative to the displayed
tail scales.  Each strict pairwise defeat of the true label implies
an error, while an error is contained in the union of weak pairwise
defeats.  Fixed \(K\), the two bounded-shift tail bounds, and fixed
ties therefore prove \cref{dcsbm:reference-sum}.
\end{proof}

\subsection{The complete degree-weighted tilted cancellation}

We now work on the genuinely deleted graph \(A^{(-i)}\).
To simplify notation within this subsection, \(D_-\), \(U_-\),
\(\Pi_-\), and
\[
 H_-=D_-^{-1/2}\Pi_-D_-^{1/2}
\]
refer only to this independent training graph.  Its completed mean
is \(\Omega_-\), and \(D_{0,-}=\diag(\Omega_-\one)\).
The tested row remains independent Bernoulli conditional on the
deleted graph.  Let \(w=(w_j)_{j\ne i}\) be the true block coefficient
vector above, and define \(e=(H_--I)w\).

\begin{lemma}[Bounded tilted coefficient error]
\label{dcsbm:weighted-cancellation}
On the deleted-graph good event, the exact profile tilt in
\cref{dcsbm:tilt} satisfies
\begin{equation}
 \max_j|e_j|\le C,\qquad
 \sum_jr_je_j^2\le C,\qquad
 \left|\sum_jr_je_j\right|\le C.
 \label{dcsbm:three-moments}
\end{equation}
All constants are uniform over the vertex and ordered pair of
distinct candidate blocks.
\end{lemma}

\begin{proof}
The residual lemma on the deleted graph gives
\[
 \|(I-\Pi_-)D_-^{1/2}Z_-\|_{\op}\le C\sqrt n.
\]
Consequently, for every bounded block vector \(w\),
\begin{equation}
 \|(I-H_-)w\|_{D_-}\le C\sqrt n,\qquad
 \|(I-H_-)w\|_2^2\le Cn/d.
 \label{dcsbm:weighted-residual}
\end{equation}
Incoherence bounds the coordinates of \(H_-w\):
\[
 |(H_-w)_j|
 \le D_{-,j}^{-1/2}\|U_{-,j\cdot}\|_2
          \|U_-^TD_-^{1/2}w\|_2
 \le (C/\sqrt d)(C/\sqrt n)C\sqrt{nd}\le C.
\]
This proves bounded coefficients.  Since \(r_j\le Cp\),
\cref{dcsbm:weighted-residual} gives
\(\sum_jr_je_j^2\le Cp\,Cn/d\le C\).

For the mean, exact Bernoulli tilting gives
\begin{equation}
 r_j=vp\theta_j C_{az_j}e^{-sw_j}+\varepsilon_j,
 \qquad \max_j|\varepsilon_j|\le Cp^2.
 \label{dcsbm:tilted-mean}
\end{equation}
For a remaining vertex in block \(c\), write
\[
 T_c^-=\sum_{\ell\ne i}\theta_\ell C_{cz_\ell},\qquad
 h_j=\frac{v C_{ac}e^{-sw_c}}{T_c^-}.
\]
Then \(h\) is block constant with \(\|h\|_\infty\le C/n\), and
the leading vector in \cref{dcsbm:tilted-mean} is exactly
\(D_{0,-}h\).  Because \(H_-\) is \(D_-\)-orthogonal,
\begin{align}
 r^T(H_--I)w
 ={}&-\langle(I-H_-)h,(I-H_-)w\rangle_{D_-}
       \nonumber\\
 &+\big((D_{0,-}-D_-)h\big)^T(H_--I)w
   +\varepsilon^T(H_--I)w .
 \label{dcsbm:weighted-identity}
\end{align}
This identity retains the random-degree correction explicitly.
The first factor in its first term is at most \(C/\sqrt n\):
apply the block residual estimate to the block coefficients of
\(h\), which have magnitude \(O(n^{-1})\).
The second factor is at most \(C\sqrt n\), so that term is \(O(1)\).

For the next term, the degree aggregate in
\cref{dcsbm:residual-estimates} gives
\[
 \|(D_{0,-}-D_-)h\|_{D_-^{-1}}^2
 =\sum_j\frac{(D_{0,-,j}-D_{-,j})^2h_j^2}{D_{-,j}}
 \le C/n.
\]
Cauchy--Schwarz in the degree metric and
\cref{dcsbm:weighted-residual} again give \(O(1)\).
Finally, bounded \(e_j\) and
\cref{dcsbm:tilted-mean} bound the last term by
\(Cnp^2=O(1)\).  This proves the tilted mean bound and the lemma.
\end{proof}

The middle term of \cref{dcsbm:weighted-identity} is why replacing
random degrees by population degrees without a remainder is invalid.
The product of two residual bounds is what reduces the leading
tilted bias to a constant.

\subsection{Completion of the risk bound on the common good event}

By \cref{dcsbm:initializer}, the initial partition lies in a fixed
\(M/d\) basin on \(\mathcal G\).  Apply one nuisance fit and one
profile assignment.  The uniform comparison
\cref{dcsbm:uniform-contrast} and the actual full-to-deleted-graph
scalar lemma \cref{lem:loo} give
\begin{equation}
 \widehat L_i^{ab}
 =\sum_{j\ne i}A_{ij}(H_-w_{i,ab,-i})_j+O_M(1)
 =T_{i,ab}+\sum_{j\ne i}A_{ij}e_j+O_M(1).
 \label{dcsbm:actual-to-loo}
\end{equation}
The leave-one-out graph appears only in this proof identity; the
implemented classifier uses the single full-graph encoding.
The \(O_M(1)\) error is a deterministic uniform bound, so it need
not be asserted to be independent of the tested row.

Condition on the genuinely deleted graph.  On its measurable good
event, \cref{dcsbm:weighted-cancellation} supplies all three
coefficient conditions of the sharp transfer lemma
\cref{lem:sharp}.  Therefore, for a fixed bounded threshold large
enough to absorb \cref{dcsbm:actual-to-loo},
\begin{equation}
 \PP_a\!\left\{
     \sum_{j\ne i}A_{ij}(H_-w_{i,ab,-i})_j\le C(M)
       \,\middle|\, A^{(-i)}
          \right\}
 \le C(M)d^{-1/2}e^{-J_{i,a\to b}} .
 \label{dcsbm:conditional-risk}
\end{equation}
The common full event implies the required minor event by the
common geometry lemmas.  Hence restricting an actual error to
\(\mathcal G\) is contained in the event on the left of
\cref{dcsbm:conditional-risk} together with the measurable minor
event.  This does not condition the tested row on a dependent
full-graph event.

Integrating over the deleted graph and summing over \(b\ne a\)
gives
\[
 \PP\{\widehat z_i\ne z_i,\mathcal G\}
 \le C\sum_{b\ne z_i}d^{-1/2}e^{-J_{i,z_i\to b}}.
\]
Here labels are aligned by the common initial permutation.
The permutation-invariant misclassification loss is no larger
than the loss under that alignment.  By
\cref{dcsbm:profile-tail}, the last sum is at most a constant
times the exact full-edge profile-reference error at \(i\).
Averaging proves
\begin{equation}
 \E\{\ell(\widehat z,z)\ind{\mathcal G}\}
 \le C R_{\mathrm{prof},*}.
 \label{dcsbm:on-good-risk}
\end{equation}
The calculation in \cref{sec:probability} proves
\(\PP(\mathcal G^c)\le C R_{\mathrm{prof},*}\), including the
minimum-degree exception at the correct exponential scale.
Because the fallback loss is at most one, adding that contribution
to \cref{dcsbm:on-good-risk} completes \cref{thm:dc}.

\subsection{Optional repeated hard fitting and finite adaptive stopping}

The same theorem holds if, after initialization, the algorithm
repeatedly recomputes \cref{dcsbm:theta-fit,dcsbm:block-criterion}
and makes a hard profile assignment using the same \(R^\circ\).
It may return any data-dependent finitely stopped iterate after
at least one completed profile assignment.  A finite iteration
cap is one explicit termination rule.

\begin{lemma}[A uniform hard-refit basin]
\label{dcsbm:basin}
There is a fixed \(M\) such that, on \(\mathcal G\) and for all
sufficiently large \(n\), every hard fit-and-profile update maps
the \(M/d\) error basin into itself, under the same label matching.
Every error at every updated iterate is contained in
\begin{equation}
 [R^\circ w_{i,ab}]_i\le C(M)
 \quad\text{for some }b\ne z_i .
 \label{dcsbm:common-error-event}
\end{equation}
\end{lemma}

\begin{proof}
First,
\begin{equation}
 \|R^\circ-\Omega\|_F^2\le C_{\rm spec}d.
 \label{dcsbm:frobenius}
\end{equation}
Indeed \(S_K=U\diag(\mu)U^T\) is a best rank-\(K\)
operator-norm approximation to \(S\).  Since \(M_D\) has rank
\(K\), \(\|S_K-M_D\|_{\op}\le2\|S-M_D\|_{\op}\le C/\sqrt d\).
Their difference has rank at most \(2K\), so its Frobenius norm
is also \(O(d^{-1/2})\).  Multiplication by \(D^{1/2}\) on both
sides gives \(\|R-\Omega\|_F\le C\sqrt d\).
The hollow term has Frobenius norm at most
\[
 \max_i|s_i|\,\|H-I\|_F\le Cp\,C\sqrt n=Cd/\sqrt n,
\]
whose square is at most \(Cd\), proving
\cref{dcsbm:frobenius}.

The completed population contrast has a uniform positive margin:
with \(\lambda_i=\sum_{j\ne i}p\theta_i\theta_jC_{z_i z_j}\),
\begin{equation}
 [\Omega w_{i,ab}]_i
 =\lambda_i\KL(q_a^{(i)}\|q_b^{(i)})+O(p)
 \ge c_{\rm gap}d,\qquad a=z_i,\ b\ne a.
 \label{dcsbm:population-margin}
\end{equation}
The \(O(p)\) term is the completed diagonal.  The uniform KL
gap follows from the same compactness and full-rank argument as
in \cref{dcsbm:profile-tail}.  Also
\(\|w_{i,ab}\|_2\le C\sqrt n\).

For an input partition in a fixed \(M/d\) basin,
\cref{dcsbm:uniform-score} bounds its fitted-contrast discrepancy
by \(C(M)\).  For large \(n\), this is at most
\(c_{\rm gap}d/2\).  A new erroneous assignment therefore requires
\[
 \|(R^\circ-\Omega)_{i\cdot}\|_2
 \ge \frac{c_{\rm gap}d}{2C\sqrt n}.
\]
Summing this squared requirement and applying
\cref{dcsbm:frobenius} shows that the new iterate has at most
\(C_{\rm new}n/d\) errors.  The constant \(C_{\rm new}\) depends
only on the fixed spectral and model constants, not on \(M\).
Choose \(M\) at least both \(C_{\rm new}\) and the initialization
constant.  All sufficiently large \(n\) then give basin invariance.
Group indices remain matched to the same true indices throughout:
the uniform score comparison and its population margin use that
fixed matching at every step.

Finally every erroneous assignment directly satisfies the
fixed-threshold canonical event in the statement by
\cref{dcsbm:uniform-score}.  The threshold is the same for all
partitions in the basin.
\end{proof}

By induction, \cref{dcsbm:basin} applies to every iterate on the
single graph event \(\mathcal G\).  The final error of an adaptively
chosen finite iterate is thus contained in the same canonical
event used to prove \cref{dcsbm:conditional-risk}.
Applying that probability bound once proves the same sharp risk
bound; there is no union over iterations and no stopping-time
independence assertion.  This corollary concerns these specified
hard updates.  It does not assert the same result for an unrelated
soft update or a reconstruction changed during iteration.

\subsection{Equal exposures and the original known-target-theta oracle}

We finish by proving \cref{cor:known}.  Suppose
\begin{equation}
 \sum_c\frac{n_c}{n}C_{ac}=S
 \quad\text{for every }a.
 \label{dcsbm:equal-exposure}
\end{equation}
The block means of theta equal one, so
\(
 T_a^{(i)}=nS-\theta_iC_{a,z_i}
\).
In particular \(T_a^{(i)}-T_b^{(i)}=O(1)\) and
\(\log(T_a^{(i)}/T_b^{(i)})=O(n^{-1})\).
Let \(v=\theta_i\) be supplied to the full-edge
known-target-theta oracle.  Its exact Bernoulli contrast is
\[
 L^{ab}_{A,\rm known}
 =\sum_cN_c\log\frac{C_{ac}}{C_{bc}}
  +\sum_{j\ne i}(1-A_{ij})
       \log\frac{1-vp\theta_jC_{az_j}}
                 {1-vp\theta_jC_{bz_j}} .
\]
On \(D_i\le Ld\), the same direct Bernoulli expansion used above
gives
\[
 L^{ab}_{A,\rm known}
 =\sum_cN_c\log\frac{C_{ac}}{C_{bc}}
   -vp(T_a^{(i)}-T_b^{(i)})+O(np^2).
\]
By \cref{dcsbm:profile-reduction}, the exact profile contrast is
\[
 L^{ab}_{A,\rm prof}
 =\sum_cN_c\log\frac{C_{ac}}{C_{bc}}
    -D_i\log\frac{T_a^{(i)}}{T_b^{(i)}}+O(np^2).
\]
Thus, using \cref{dcsbm:equal-exposure},
\begin{equation}
 |L^{ab}_{A,\rm prof}-L^{ab}_{A,\rm known}|
 \le C(np^2+p)=O(1).
 \label{dcsbm:known-comparison}
\end{equation}
Fixed prior log-ratios, if present, are identical on both sides.
The bounded-shift upper and lower tails in
\cref{dcsbm:directional-tail} show that the two directional
error probabilities are comparable by constants.
The upper-degree exception has arbitrarily large fixed
exponential rate after enlarging \(L\), and is negligible for
both references.  Finite unions over competitors and averaging
over vertices yield
\[
 R_{\mathrm{prof},*}\asymp R_{\theta,*}.
\]
Combining this with \cref{thm:dc} proves \cref{cor:known}.
The condition permits heterogeneous theta, unequal block sizes,
and general positive full-rank symmetric \(C\) obeying its exposure
constraints.  Its role is precisely to remove the label information
carried by the leading total-degree exposure, which profiling
otherwise changes.

\subsection{The classical balanced DC-SBM oracle and its sharp risk}
\label{dcsbm:classical-oracle}

The fundamental local testing problem used by
Gao, Ma, Zhang and Zhou~\cite{gmzz2018} fixes the target degree parameter
under both candidate labels. Its exact simple-hypothesis likelihood test
therefore corresponds to \(R_{\theta,*}\) in \cref{sec:oracle}.
The next proposition explicitly connects that reference, the profile
reference and true-group neighbor voting in the balanced planted-partition
model. This is an oracle comparison with its own assumptions:
it uses neither the spectral good event \(\mathcal G\) nor a logarithmic
lower bound on \(np\).

\begin{proposition}[Sharp alignment of the classical DC-SBM references]
\label{dcsbm:classical-bridge}
\label{prop:classical-dc-oracle-bridge}
Fix an arbitrary \(K\ge2\), positive bounds
\(\theta_-,\theta_+\), and \(0<r_-<r_+<1\).
Let \(n=Km\), let every group have \(m\) vertices, and suppose
\[
 \theta_-\le\theta_i\le\theta_+,\qquad
 \sum_{i:z_i=a}\theta_i=m\quad\text{for all }a.
\]
Specialize the block matrix \(B=pC\) to
\begin{equation}
 B_{aa}=p,\qquad B_{ab}=q\ (a\ne b),\qquad
 r_-\le q/p\le r_+ .
 \label{dcsbm:classical-model}
\end{equation}
Assume only, for the present oracle comparison,
\begin{equation}
 d=np\longrightarrow\infty,\qquad np^2\le C_1<\infty.
 \label{dcsbm:classical-oracle-range}
\end{equation}
Define \(R_{\mathrm{vote},*}\) as the average error of the local raw-row
rule which knows the neighbor labels and selects the group with the
largest incident block count. Use any fixed tie ordering.
Then, uniformly over this parameter class,
\begin{equation}
 \begin{aligned}
 R_{\theta,*}
 &\asymp R_{\mathrm{prof},*}
 \asymp R_{\mathrm{vote},*}\\
 &\asymp
 \frac1n\sum_{i=1}^n
 \frac1{\sqrt{np}}
 \exp\left\{-\theta_i\frac nK(\sqrt p-\sqrt q)^2\right\}.
 \end{aligned}
 \label{dcsbm:classical-risk}
\end{equation}
The profile reference is exactly the Bernoulli profile in
\eqref{eq:profile-oracle}, with candidate probabilities capped at \(1/2\).
The same comparison holds if the voting counts are normalized by the
self-excluded neighbor-group sizes. Constants depend only on the displayed
fixed bounds, not on \(n\). In particular,
\eqref{dcsbm:classical-oracle-range} permits \(np=o(\log n)\).
\end{proposition}

\begin{proof}
Fix a target \(i\), its true group \(a\), and a competitor \(b\ne a\).
For this proof write
\[
 N_c=\sum_{\substack{j\ne i\\z_j=c}}A_{ij},\qquad
 M_c=\sum_{\substack{j\ne i\\z_j=c}}\theta_j
     =m-\theta_i\ind{c=a},\qquad
 L=\log(p/q).
\]
Thus \(L\) is bounded above and below by positive constants.
Put
\[
 V_{i,ab}=L(N_a-N_b),\qquad
 J_i=\theta_i m(\sqrt p-\sqrt q)^2 .
\]
The assumptions imply \(J_i\asymp d\), uniformly in \(i\).
They also imply \(p\to0\), so every true candidate edge probability
is eventually less than \(1/2\).

\medskip
\noindent\textbf{The exact Bernoulli count transform.}
Under the actual true-row law, independence gives
\begin{align}
 \Psi_{i,ab}(s)
 &:=\log\E_a e^{-sV_{i,ab}}\nonumber\\
 &=\sum_{\substack{j\ne i\\z_j=a}}
   \log\{1+\theta_i\theta_jp(e^{-sL}-1)\}
   +\sum_{j:z_j=b}
   \log\{1+\theta_i\theta_jq(e^{sL}-1)\}.
 \label{dcsbm:classical-cumulant}
\end{align}
Taylor expansion, uniformly for \(s\in[0,1]\), gives
\begin{equation}
 \Psi_{i,ab}(s)
 =m\theta_i\{p(e^{-sL}-1)+q(e^{sL}-1)\}
       +O(np^2+p).
 \label{dcsbm:classical-cumulant-expansion}
\end{equation}
The same error bound holds for the first two derivatives. Bounded theta
and bounded \(L\) make the sum of quadratic Bernoulli remainders
\(O(np^2)\). The additional \(O(p)\) term accounts for the omitted
\(\theta_i\) in \(M_a=m-\theta_i\).

The leading transform in
\eqref{dcsbm:classical-cumulant-expansion} is minimized at \(s=1/2\),
with minimum \(-J_i\). Its endpoint derivatives have opposite signs,
each of magnitude at least \(cd\), and its second derivative is
between positive constant multiples of \(d\).
Since \(np^2+p=o(d)\), the exact minimizing saddle stays in a fixed
compact subinterval of \((0,1)\). Convexity makes it the global
minimum over \(s\ge0\). In particular,
\begin{equation}
 I^{\mathrm{cnt}}_{i,ab}
 :=-\min_{s\ge0}\Psi_{i,ab}(s)
 =J_i+O(np^2+p)=J_i+O(1).
 \label{dcsbm:classical-information}
\end{equation}
One can also obtain this bound directly by taking minima in the
uniform transform approximation.

Under the exact optimizing tilt, the Bernoulli probabilities on the
two relevant groups are comparable to \(p\). Their sum is \(O(d)\),
and the variance of \(V_{i,ab}\) is comparable to \(d\).
Consequently \cref{lem:sharp} applies to
\eqref{dcsbm:classical-cumulant} without any spectral event. For every
fixed \(x_0<\infty\), uniformly over \(|x|\le x_0\),
\begin{equation}
 c_{x_0}d^{-1/2}e^{-J_i}
 \le \PP_a(V_{i,ab}<x)
 \le \PP_a(V_{i,ab}\le x)
 \le C_{x_0}d^{-1/2}e^{-J_i}.
 \label{dcsbm:classical-count-tail}
\end{equation}
The strict lower bound from that lemma uses an open interval of fixed
length under the tilt. Thus lattice atoms and an arbitrary fixed
tie ordering cause no loss of the lower bound.

\medskip
\noindent\textbf{Comparison with the known-target likelihood.}
Let
\[
 \mathcal T_c=\sum_{j\ne i}\theta_j B_{c,z_j}
              =pT_c^{(i)}.
\]
By the exact group normalization,
\begin{equation}
 \begin{aligned}
 \mathcal T_a&=m[p+(K-1)q]-\theta_ip,\\
 \mathcal T_b&=m[p+(K-1)q]-\theta_iq.
 \end{aligned}
 \label{dcsbm:classical-exposure}
\end{equation}
In particular \(\mathcal T_a,\mathcal T_b\asymp d\),
\(\mathcal T_a-\mathcal T_b=O(p)\), and
\(\log(\mathcal T_a/\mathcal T_b)=O(n^{-1})\).

Fix a sufficiently large constant \(L_0\) to be chosen below, and consider
the single-target event
\[
 \mathcal E_i=\{D_i\le L_0d\}.
\]
The exact known-target Bernoulli contrast is
\[
 L^{ab}_{i,\mathrm{known}}
 =V_{i,ab}+
 \sum_{j\ne i}(1-A_{ij})
 \log\frac{1-\theta_i\theta_jB_{a,z_j}}
          {1-\theta_i\theta_jB_{b,z_j}} .
\]
Expansion of these nonedge logarithms on \(\mathcal E_i\) gives
\begin{align}
 L^{ab}_{i,\mathrm{known}}
 &=V_{i,ab}-\theta_i(\mathcal T_a-\mathcal T_b)
       +O(D_ip+np^2)\nonumber\\
 &=V_{i,ab}+O(np^2+p).
 \label{dcsbm:classical-known-contrast}
\end{align}
The \(O(D_ip)\) term restores the incident-edge terms in the linear
exposure; the quadratic remainder is \(O(np^2)\).

\medskip
\noindent\textbf{Comparison with the capped Bernoulli profile.}
For a candidate \(c\), put \(b_{cj}=\theta_jB_{c,z_j}\).
The admissible interval in \eqref{eq:profile-oracle} is
\[
 0\le t\le t_{\max,c}:=\frac1{2\max_{j\ne i}b_{cj}}.
\]
Because \(b_{cj}\asymp p\), its upper endpoint tends to infinity.
For \(t>0\) the derivative of the exact raw-row likelihood is
\[
 \frac{D_i}{t}
 -
 \sum_{j\ne i}(1-A_{ij})
       \frac{b_{cj}}{1-tb_{cj}} .
\]
On \(\mathcal E_i\), there are at least \(n-1-L_0d\ge n/2\)
nonedges for all sufficiently large \(n\). Their total intensity is
at least \(cd\), uniformly over candidates. Hence the derivative is
negative for every admissible \(t\) above a fixed constant depending
only on \(L_0\) and the structural bounds. Every exact maximizing
\(t\) is therefore bounded. The cap at \(1/2\) is inactive on this
bounded interval for all large \(n\).

Uniformly over the resulting bounded \(t\) interval, the likelihood is
\begin{equation}
 D_i\log t+\sum_{j\ne i}A_{ij}\log b_{cj}
      -t\mathcal T_c+O(np^2).
 \label{dcsbm:classical-profile-expansion}
\end{equation}
Indeed, restoring the edge terms in the linear exposure costs
\(O(D_ip)\), and the sum of quadratic nonedge remainders costs
\(O(np^2)\). Both are \(O(np^2)\) on \(\mathcal E_i\).
The leading optimizer \(D_i/\mathcal T_c\) is bounded as well, so
the same error bound holds after maximization.
Terms \(D_i\log D_i-D_i\) and
\(\sum_{j\ne i}A_{ij}\log\theta_j\) cancel in a contrast.
Thus the exact capped profile satisfies
\begin{align}
 L^{ab}_{i,\mathrm{prof}}
 &=V_{i,ab}
   -D_i\log(\mathcal T_a/\mathcal T_b)+O(np^2)\nonumber\\
 &=V_{i,ab}+O(np^2+p)
 \quad\text{on }\mathcal E_i.
 \label{dcsbm:classical-profile-contrast}
\end{align}
For \(D_i=0\), the raw row is identically zero; \(t=0\) maximizes
every candidate likelihood with value zero. The displayed formulas
then hold with the continuous convention \(0\log0=0\).

\medskip
\noindent\textbf{Single-row exceptions, directional errors, and averaging.}
Every mean degree is at most \(\lambda_+d\) for a fixed \(\lambda_+\).
For any prescribed fixed \(H\), a Bernoulli Chernoff bound permits a
fixed choice of \(L_0\) such that
\begin{equation}
 \sup_i\PP(\mathcal E_i^c)\le e^{-Hd}.
 \label{dcsbm:classical-single-exception}
\end{equation}
Also \(J_i\le C_Jd\) uniformly. Choose \(H>C_J+1\).
It follows from \(d\to\infty\) that
\[
 \sup_i
 \frac{\PP(\mathcal E_i^c)}
      {d^{-1/2}e^{-J_i}}
 \le \sqrt d\,e^{-(H-C_J)d}\longrightarrow0 .
\]
This calculation concerns one target at a time; it introduces no
factor \(n\) and requires no logarithmic lower bound on \(d\).

By \eqref{dcsbm:classical-known-contrast} and
\eqref{dcsbm:classical-profile-contrast}, a fixed constant \(B_0\)
bounds the difference between \(V_{i,ab}\) and either exact contrast
on \(\mathcal E_i\). Thus for either
\(L^{ab}_i=L^{ab}_{i,\mathrm{known}}\) or
\(L^{ab}_i=L^{ab}_{i,\mathrm{prof}}\),
\begin{align*}
 \PP_a(V_{i,ab}<-B_0)-\PP_a(\mathcal E_i^c)
 &\le \PP_a(L^{ab}_i<0)\\
 &\le \PP_a(L^{ab}_i\le0)\\
 &\le \PP_a(V_{i,ab}\le B_0)+\PP_a(\mathcal E_i^c).
\end{align*}
Applying \eqref{dcsbm:classical-count-tail} and then
\eqref{dcsbm:classical-single-exception} proves that both the strict
and nonstrict directional errors of either oracle are comparable
to \(d^{-1/2}e^{-J_i}\).
These are unconditional event inclusions; no independence from
\(\mathcal E_i\) is asserted or needed.

For the raw block-count vote, the pairwise contrast is \(V_{i,ab}\)
itself. If counts are normalized by the self-excluded group sizes,
multiplying all candidate scores by \(mL\) changes the pairwise
contrast by at most \(CLD_i/n=O(p)\) on \(\mathcal E_i\), since
those sizes are \(m\) or \(m-1\). The same tail comparison therefore
applies to either voting convention.

For each of these coherent families of candidate scores, a strict
pairwise defeat of the true group implies an error, and an error
implies a nonstrict defeat by some competitor. Taking the maximum
of the strict pair probabilities for a lower bound and their
nonstrict sum for an upper bound proves that each per-target error
is comparable to \(d^{-1/2}e^{-J_i}\), since \(K\) is fixed.
Averaging these bounds over \(i\) proves
\eqref{dcsbm:classical-risk}.
\end{proof}

\begin{remark}[Oracle alignment and the algorithmic range]
\label{dcsbm:classical-scope}
The degree-dependent exponent in \eqref{dcsbm:classical-risk} is the
balanced DC-SBM exponent appearing in the minimax analysis of
Gao, Ma, Zhang and Zhou~\cite{gmzz2018}. Their exponent-level
minimax expression does not, by itself, give the sharp prefactor
or a constant-ratio theorem for every parameter point.
The proposition supplies the exact local-oracle comparison on the
fixed-separation, bounded-theta class stated here.

For \eqref{dcsbm:classical-model}, the finite-\(n\) exposures in
\eqref{eq:equal-exposure} are exactly equal. Hence, on the overlap
with the algorithmic range of \cref{thm:dc},
\cref{cor:known,dcsbm:classical-bridge} give
\[
 \E\ell(\widehat z_D,z)\le C R_{\theta,*},
 \qquad R_{\theta,*}\asymp R_{\mathrm{prof},*}.
\]
The absence of a logarithmic lower bound in the oracle proposition
does not remove that bound from the spectral algorithm theorem.
Extending the algorithm requires its own sparse spectral and
likelihood-transfer estimates.
\end{remark}

